\documentclass[../BTL.tex]{subfiles}
\begin{document}

\section{Thiết kế cơ sở dữ liệu}
Hệ thống sử dụng cơ sở dữ liệu quan hệ với 8 bảng chính được thiết kế
chặt chẽ. Sơ đồ thực thể liên kết (ERD) và đặc tả cấu trúc các bảng được
mô tả chi tiết như sau:

\textbf{1. Bảng users (Tài khoản người dùng)} Lưu trữ thông tin của
khách hàng, quản trị viên và dữ liệu gốc của thợ.

\begin{itemize}
\item
  user\_id (INT, Primary Key): Mã người dùng.
\item
  full\_name (VARCHAR): Họ và tên đầy đủ.
\item
  email (VARCHAR): Địa chỉ email.
\item
  phone (VARCHAR): Số điện thoại đăng nhập.
\item
  password (VARCHAR): Mật khẩu.
\item
  role \textbf{(ENUM(\textquotesingle customer\textquotesingle,
  \textquotesingle stylist\textquotesingle,
  \textquotesingle admin\textquotesingle))}: Vai trò phân quyền tài
  khoản (Khách hàng, Thợ cắt tóc, hoặc Quản trị viên).
\item
  created\_at (TIMESTAMP): Thời gian tạo tài khoản.
\end{itemize}

\textbf{2. Bảng branches (Chi nhánh)} Lưu trữ thông tin các cơ sở của
Barbershop.

\begin{itemize}
\item
  branch\_id (INT, Primary Key): Mã chi nhánh.
\item
  branch\_name (VARCHAR): Tên chi nhánh.
\item
  address (VARCHAR): Địa chỉ chi nhánh.
\item
  phone (VARCHAR): Số điện thoại cơ sở.
\item
  email (VARCHAR): Email liên hệ cơ sở.
\end{itemize}

\textbf{3. Bảng stylists (Đội ngũ thợ cắt tóc)} Lưu trữ hồ sơ làm việc
của thợ, được liên kết từ bảng người dùng.

\begin{itemize}
\item
  stylist\_id (INT, Primary Key): Mã hồ sơ thợ.
\item
  user\_id (INT, Foreign Key): Liên kết với bảng users.
\item
  branch\_id (INT, Foreign Key): Trực thuộc chi nhánh nào (Liên kết bảng
  branches).
\item
  status \textbf{(ENUM(\textquotesingle active\textquotesingle,
  \textquotesingle inactive\textquotesingle))}: Trạng thái làm việc
  (Đang hoạt động hoặc Đã khóa/Nghỉ việc).
\end{itemize}

\textbf{4. Bảng services (Dịch vụ)} Lưu trữ danh sách các dịch vụ cắt
tóc.

\begin{itemize}
\item
  service\_id (INT, Primary Key): Mã dịch vụ.
\item
  service\_name (VARCHAR): Tên dịch vụ.
\item
  price (INT): Giá tiền dịch vụ.
\item
  duration (INT): Thời gian thực hiện (phút).
\item
  description (TEXT): Mô tả chi tiết dịch vụ.
\end{itemize}

\textbf{5. Bảng appointments (Lịch hẹn)} Lưu trữ chi tiết các đơn đặt
lịch của khách hàng.

\begin{itemize}
\item
  appointment\_id (INT, Primary Key): Mã lịch hẹn.
\item
  customer\_id (INT, Foreign Key): Mã khách đặt (Liên kết bảng users).
\item
  branch\_id (INT, Foreign Key): Mã chi nhánh (Liên kết bảng branches).
\item
  service\_id (INT, Foreign Key): Mã dịch vụ (Liên kết bảng services).
\item
  stylist\_id (INT, Foreign Key): Mã thợ thực hiện (Liên kết bảng
  stylists).
\item
  appointment\_date (DATE): Ngày đặt lịch.
\item
  appointment\_time (TIME): Giờ đặt lịch.
\item
  notes (TEXT): Ghi chú của khách hàng.
\item
  total\_price (INT): Tổng tiền hóa đơn.
\item
  status \textbf{(ENUM(\textquotesingle pending\textquotesingle,
  \textquotesingle confirmed\textquotesingle,
  \textquotesingle completed\textquotesingle,
  \textquotesingle cancelled\textquotesingle))}: Trạng thái lịch hẹn
  (Chờ duyệt, Đã xác nhận, Hoàn thành, hoặc Đã hủy).
\item
  created\_at (TIMESTAMP): Thời gian tạo đơn.
\end{itemize}

Trong nghiệp vụ đặt lịch, hệ thống còn hỗ trợ lọc danh sách thợ VIP theo
lựa chọn của khách hàng và kiểm tra khung giờ trống trước khi xác nhận
lịch hẹn. Cách xử lý này giúp hạn chế trùng lịch, đồng thời đảm bảo
người dùng chỉ có thể chọn các khung giờ thực sự còn khả dụng.

\textbf{6. Bảng stylist\_schedules (Lịch làm việc của thợ)} Kiểm soát
khung thời gian thực tế của thợ để tránh trùng lịch.

\begin{itemize}
\item
  schedule\_id (INT, Primary Key): Mã lịch trình.
\item
  stylist\_id (INT, Foreign Key): Mã thợ (Liên kết bảng stylists).
\item
  appointment\_id (INT, Foreign Key, Nullable): Mã lịch hẹn tương ứng
  (Liên kết bảng appointments).
\item
  work\_date (DATE): Ngày làm việc.
\item
  start\_time (TIME): Giờ bắt đầu.
\item
  end\_time (TIME): Giờ kết thúc.
\item
  status \textbf{(ENUM(\textquotesingle available\textquotesingle,
  \textquotesingle booked\textquotesingle))}: Trạng thái của khung giờ
  (Còn trống hoặc Đã được đặt).
\end{itemize}

\textbf{7. Bảng leave\_requests (Đơn xin nghỉ phép)} Lưu trữ yêu cầu xin
nghỉ của thợ để Admin phê duyệt.

\begin{itemize}
\item
  request\_id (INT, Primary Key): Mã đơn xin nghỉ.
\item
  stylist\_id (INT, Foreign Key): Mã thợ (Liên kết bảng stylists).
\item
  start\_date (DATE): Ngày bắt đầu nghỉ.
\item
  end\_date (DATE): Ngày kết thúc nghỉ.
\item
  leave\_type (VARCHAR): Loại nghỉ phép.
\item
  reason (TEXT): Lý do xin nghỉ.
\item
  status \textbf{(ENUM(\textquotesingle pending\textquotesingle,
  \textquotesingle approved\textquotesingle,
  \textquotesingle rejected\textquotesingle))}: Trạng thái đơn (Chờ
  duyệt, Đã duyệt, hoặc Bị từ chối).
\item
  created\_at (TIMESTAMP): Thời gian tạo đơn.
\end{itemize}

\textbf{8. Bảng reviews (Đánh giá dịch vụ)} Lưu trữ thông tin phản hồi,
điểm số đánh giá và bình luận của khách hàng sau khi hoàn thành lịch hẹn
cắt tóc.

\begin{itemize}
\item
  review\_id (INT, Primary Key): Mã định danh độc nhất cho mỗi lượt đánh
  giá (Tự động tăng).
\item
  appointment\_id (INT, Foreign Key, Unique): Liên kết với bảng
  appointments. Đảm bảo ràng buộc mỗi lịch hẹn chỉ được đánh giá duy
  nhất một lần để tránh spam dữ liệu.
\item
  customer\_id (INT, Foreign Key): Mã khách hàng thực hiện đánh giá
  (Liên kết với bảng users).
\item
  stylist\_id (INT, Foreign Key): Mã thợ cắt tóc nhận đánh giá (Liên
  kết với bảng users).
\item
  rating (INT): Số sao đánh giá từ khách hàng. Ràng buộc điều kiện
  (CHECK) giá trị nghiêm ngặt trong khoảng từ 1 đến 5 sao.
\item
  comment (TEXT): Nội dung bình luận chi tiết từ khách hàng (Có thể để
  trống).
\item
  created\_at (TIMESTAMP): Thời gian hệ thống ghi nhận lượt đánh giá
  ban đầu.
\item
  updated\_at (TIMESTAMP): Thời gian cập nhật nếu khách hàng có sự sửa
  đổi nội dung đánh giá.
\end{itemize}

\placedfigure{Hinhve/image34.png}{Sơ đồ thực thể liên kết (ERD)}

\section{Tổng kết chương}
Chương 3 đã trình bày tổng quan về usecase tổng quan hệ thống, usecase
chi tiết cùng với kịch bản usecase, các chức năng chi tiết trong hệ
thống với từng vai trò, bên canh đó còn thể hiện biểu đồ tuần tự, cơ sở
dữ liệu.

\protect\hypertarget{_Toc225990067}{}{}\textbf{\hfill\break
}

\end{document}
